% === 计算数学 (Journal of Computational Mathematics) === 中文数学期刊排版模板（YangTex 内置）
% 引擎：xelatex + ctex；按期刊投稿规范整理，非期刊官方 .cls 文件。
% 分类：中文数学期刊
\documentclass[UTF8,fontset=fandol,11pt]{ctexart}
\usepackage{amsmath,amssymb,amsthm,mathtools}
\usepackage{graphicx}
\newtheorem{theorem}{定理}[section]
\newtheorem{lemma}[theorem]{引理}
\newcommand{\R}{\mathbb{R}}
\newcommand{\Om}{\Omega}
\newcommand{\dd}{\,\mathrm{d}}
\title{非局部椭圆问题有限元方法的最优收敛阶}
\author{杨建伟}
\date{}
\begin{document}
\maketitle
\begin{abstract}
针对非局部椭圆算子构造协调有限元离散格式，在最低正则性假设下证明了能量范数下的最优阶误差估计，并用数值实验验证。
\end{abstract}
\paragraph{关键词} 有限元方法；非局部算子；先验误差估计；收敛阶

\section{引言}
在形状正则剖分 $\mathcal{T}_h$ 上用 Galerkin 方法离散积分算子 $(-\Delta)^s$。
经典背景与相关结果可参见~\cite{Evans,Brezis}。

\begin{theorem}
设真解 $u\in H^{1+s}(\Om)$，则离散解满足 $\|u-u_h\|_{H^s}\le Ch^{1-s}|u|_{H^{1+s}}$。
\end{theorem}

\begin{proof}
由 Galerkin 正交性、插值估计与离散双线性型的稳定性即得结论。
\end{proof}

\section{主要估计}
\begin{equation}
a(u_h,v_h)=\int_\Om f v_h\,\dd x\qquad\forall v_h\in V_h,
\end{equation}
均匀网格与分级网格上的数值结果均与理论阶一致。

\section{结论}
本文的结果可推广到更一般的区域与边界条件，相关讨论将在后续工作中给出。

\begin{thebibliography}{9}
\bibitem{Evans} 伍卓群、尹景学、王春朋，\emph{椭圆与抛物型方程引论}, 科学出版社, 2013.
\bibitem{Brezis} H. Brezis, \emph{Functional Analysis, Sobolev Spaces and PDEs}, Springer, 2011.
\end{thebibliography}
\end{document}
